\section{结论}
\label{sec:conclusion}

\policy{} 从各问题需要不同迭代次数的批次中移除 padding OT 工作。该方法以独立 Sinkhorn 问题的标准固定宽度执行为起点；所选求解器的停止规则判断问题是否完成，动态压缩则使此后的逐问题 kernel 变窄。生存曲线恒等式计算消失的问题更新数。机会--转化分解随后将存活宽度映射到实测后端成本曲线并减去控制开销，从而给出匹配路径的精确 break-even 条件。

匹配的四 worker MetroPT-3 端点达到 \CMetroCorrectedEndpointSpeedup{}，ImageNet-32 训练 solver 字段达到 $1.054\times$，已登记的 MetroPT scale cell 达到 $1.616$--$3.110\times$。多个问题使用一个初始化的非冷启动条件达到 $1.421\times$，并减少 30.62\% 的问题更新槽位。独立的 OTT-JAX 和 PyKeOps 实现通过各自后端移除已完成问题。

匹配对照识别了加速机制。批量检查与 Sinkhorn 专用预检降低检查成本；逻辑 masking 保持 kernel 宽度不变；动态压缩移除 OT 工作。在已登记的 MetroPT 与 Packer19 端点窗口中，动态压缩同时降低实测 GPU 能耗。在保持总迭代次数不变的条件下进行分组，结果说明窗口内所需迭代编号的差异越大，可移除 padding 越多。直接 A100 宽度扫描呈现波次形状的转化曲线：较小宽度可能共享同一成本层级，而更大的问题会更早地响应宽度缩减。报告的应用把速度与实际残差以及消费者或训练端点配对呈现。快速内部 kernel 让每次更新更便宜；DrainSinkhorn 则不再向已完成问题发出这些更新。

\paragraph{可复现性。}
发布制品记录了源代码和输入身份、命令、软件与 GPU 环境、原始输出、数值检查和配对分析。
